\newif\ifglobal

%%%%%%%%%%%%%%%%%%%%%%%
%%% Compile to view %%%
%%%%%%%%%%%%%%%%%%%%%%%

%%%%%%%%%%%%%%%%%%%%%%%%%%%%%%%%%%%%%%%%%%%%%%%%%%%
%%% Readme:                                     %%%
%%% https://www.overleaf.com/read/bwdjvfjksvjy/ %%%
%%%%%%%%%%%%%%%%%%%%%%%%%%%%%%%%%%%%%%%%%%%%%%%%%%%

% >>> M?: marks a module setting number or important notice

% >>> M4: Set {./relative-path} to external files for this module <<<
\def \modulefiles {./22-RNG}
% To add an external file, use:
% (Replace '---' with actual filename)
% '\input{\modulefiles/tables/---.tex}' for tables
% '\includegraphics{\modulefiles/figures/---}' for figures
% '\subfile{\modulefiles/---}' for register files


% >>> M1: This module is in global project? <<<
% 'YES' - uncomment; 'NO' - comment out
\globaltrue

\ifglobal
    \documentclass[main\_\_CN.tex]{subfiles}

\else
    % Font "FandolSong-Regular" used by ctex does not have GSUB table.
    % It causes warnings that do not affect compilation and can be ignored.
    \PassOptionsToPackage{quiet}{fontspec}

    \documentclass[openany, 10pt]{book}

    % >>> M2: In ./chip-spec-settings/preamble-trm-module, also find
    %         settings for visibility of todo notes, watermark <<<
    \usepackage{preamble-trm-module}


    % Enable Chinese typesetting
    \usepackage{ctex}

    % >>> M3: Temporary labels in Chinese
    %         you can add your own labels there <<<
    \myexternaldocument{temp-labels-cn}

\fi %\ifglobal


\setcounter{secnumdepth}{4}



% >>> M5: If this module needs any updates in
% './chip-spec-settings/preamble-trm-module.sty'
% (include additional package, etc.), contact your module PM
% or create the file './chip-spec-settings/changelog.tex' make the
% necessary updates yourself and document them in this file <<<

\raggedright
\raggedbottom


\begin{document}


% Sets language into which pre-defined bits of text are translated
\selectlanguage{chinese}
\setchineseformatting

% Do not delete this block
\ifglobal
% Add the following front matter if
% this module is outside of global project
\else
    \InputIfFileExists{readme}{}{}
    {\let\clearpage\relax \listoftodos}
    {\let\clearpage\relax \listofdonetodos}
    \tableofcontents
    %\listoftables
    %\listoffigures
\fi


%%%%%%%%%%%%%%%%%%%%%%%%%
%%% TRM Module Begins %%%
%%%%%%%%%%%%%%%%%%%%%%%%%

\hypertarget{rng}{}
\chapter{随机数发生器 (RNG)}\label{mod:rng}

\section{概述}

ESP32-S2 内置一个真随机数发生器，其生成的 32 位随机数可作为加密等操作的基础。

\section{主要特性}

ESP32-S2 的随机数发生器可通过物理过程而非算法生成真随机数，所有生成的随机数在特定范围内出现的概率完全一致。

\section{功能描述}\label{sec:rng-func-descr}

系统可以从随机数发生器的寄存器 \hyperref[regdesc:RNGDATAREG]{RNG\_DATA\_REG} 中读取随机数，每个读到的 32 位随机数都是真随机数，噪声源为系统中的热噪声和异步时钟。

具体来说，这些热噪声可以来自 SAR ADC 或高速 ADC 或两者兼有。当芯片的 SAR ADC 或高速 ADC 工作时，就会产生比特流，并通过异或 (XOR) 逻辑运算作为随机数种子进入随机数生成器。

当为数字内核使能 RC\_FAST\_CLK 时钟时，随机数发生器也会对 RC\_FAST\_CLK (8 MHz) 进行采样，作为随机数种子。RC\_FAST\_CLK 是一种异步时钟源，由于存在电路亚稳态，因此可以提高随机数发生器的墒值。然而，为了保证随机数发生器可以获得最大熵值，仍建议在使用随机数发生器时至少使能一路 ADC 作为随机数种子。

\begin{figure}[H]
    \centering
    \includegraphics[width=0.9\textwidth]{\modulefiles/figures/06-RNG-noice-source}
    \caption{噪声源}
    \label{figure:RNG_Noice_Source}
\end{figure}


当 SAR ADC 打开时，每个 RC\_FAST\_CLK （8 MHz）时钟周期内（来自内部 RC 振荡器，详见 \ref{mod:resclk} \textit{\nameref{mod:resclk}} 章节），随机数发生器将获得 2 位的熵。因此，为了获得最大的熵值，建议读取 \hyperref[regdesc:RNGDATAREG]{RNG\_DATA\_REG} 寄存器时的速率不超过 500 kHz。

当高速 ADC 打开时，每个 APB 时钟周期（通常为 80 MHz）内，随机数发生器将获得 2 位的熵。因此，为了获得最大的熵值，建议读取 \hyperref[regdesc:RNGDATAREG]{RNG\_DATA\_REG} 寄存器时的速率不超过 5 MHz。



\section{编程指南}\label{rng-programming}

在使用 ESP32-S2 的随机数生成器时，应该至少打开 SAR ADC，高速 ADC，或 RC\_FAST\_CLK，否则可能会导致产生伪随机数，应注意避免。其中，
\begin{itemize}
    \item SAR ADC 可通过 DIG ADC 控制器打开，详见 \ref{mod:sensor} \textit{\nameref{mod:sensor}} 章节。
    \item 高速 ADC 在 Wi-Fi 开启时自动打开。
    \item RC\_FAST\_CLK 可通过设置 \hyperref[regdesc:RTCCNTLCLKCONFREG]{RTC\_CNTL\_CLK\_CONF\_REG} 寄存器中的 \hyperref[fielddesc:RTCCNTLDIGCLK8MEN]{RTC\_CNTL\_DIG\_CLK8M\_EN} 位使能。
\end{itemize}

\begin{tiplisting}
注意，在 Wi-Fi 开启时，极端情况下高速 ADC 有读值饱和的可能，这会降低熵值。因此，建议在 Wi-Fi 开启时，同时通过 DIG ADC1 控制器打开 SAR ADC 产生随机数。
\end{tiplisting}

在使用随机数生成器时，请多次读取 \hyperref[regdesc:RNGDATAREG]{RNG\_DATA\_REG} 寄存器的值，直至获得足够多的随机数。在读取寄存器时，注意控制速率不要超过上方第 \ref{sec:rng-func-descr} 小节的介绍。


\section{基地址}\label{subsection:RNG-base-address}

用户可以通过两个不同的寄存器基地址访问随机数发生器，如表 \ref{table:RNG-base-address} 所示。更多信息，请访问章节 \ref{mod:sysmem} \textit{\nameref{mod:sysmem}}。

\begin{longtable}[c]{ | l | C{3cm} | }
\caption{随机数发生器基地址}\label{table:RNG-base-address}
\\\hline
\rowcolor{lightgray}
    访问总线 & 基地址\\
    \hline
    PeriBUS1      & \hexprefix{}3F435000\\
    \hline
    PeriBUS2      & \hexprefix{}60035000\\
    \hline
\end{longtable}

\hypertarget{rng-reg-summ}{}
\section{寄存器列表}

请注意，下表中的地址都是相对于随机数发生器基地址的地址偏移量（相对地址）。更多有关随机数发生器基地址的信息，请前往 \ref{subsection:RNG-base-address} 章节。
\begin{longtable}{ | p{5cm} | p{3cm} | p{3cm} | p{1cm} | }
\hline
\rowcolor{lightgray}
名称 & 描述 & 地址 & 访问 \\
\hline
\hyperref[regdesc:RNGDATAREG]{RNG\_DATA\_REG} & 随机数数据 & \hexprefix{}0110 & 只读 \\
\hline
\end{longtable}


\section{寄存器}
请注意，这里的地址都是相对于随机数发生器基地址的地址偏移量（相对地址）。更多有关随机数发生器基地址的信息，请前往 \ref{subsection:RNG-base-address} 章节。
\begin{register}{H}{RNG\_DATA\_REG}{\hexprefix{}0110}\label{regdesc:RNGDATAREG}
\regfield{}{32}{0}{{\hexprefix{}00000000}}%
\reglabel{Reset}\regnewline%
\begin{regdesc}\begin{reglist}
\item [RNG\_DATA]随机数来源。（只读）
\end{reglist}\end{regdesc}
\end{register}


\end{document}
